% !TeX program = xelatex
% XA-202620 面向政企场景的大模型智能体安全关键技术研究 技术方案报告
%
% 编译: xelatex main.tex (两遍, 或用 latexmk -xelatex main.tex)
% 依赖: TeX Live 2020+, ctex 宏集
%
% 写作规范(全体作者遵守):
%   1. 不使用破折号(——), 用逗号、句号或"即"替代
%   2. 不使用"不是…而是"句式
%   3. 不做防御性写作, 不预答辩, 不自我贬低
%   4. 不提及任何借鉴来源, 全文以自主研制口径陈述
%   5. 数字优先, 能用数字不用形容词
%   6. 每段首句说清本段要干什么
%
% 结构依据: docs/competition/章节规划-v2.md
% 官方要求: 问题分析/技术路线/算法设计/实验方案/指标体系/预期效果, ≤30 页

\documentclass[UTF8, a4paper, 12pt]{ctexart}

\usepackage[margin=2.5cm]{geometry}
\usepackage{booktabs}
\usepackage{graphicx}
\usepackage{float}
\usepackage{listings}
\usepackage{xcolor}
\usepackage{array}
\usepackage{longtable}
\usepackage{enumitem}
\usepackage{caption}
\usepackage{hyperref}

\hypersetup{colorlinks=true, linkcolor=black, urlcolor=blue, citecolor=black}
\captionsetup{font=small, labelsep=quad}

\lstset{
  basicstyle=\ttfamily\small,
  breaklines=true,
  frame=single,
  backgroundcolor=\color{gray!8},
  columns=flexible,
  keepspaces=true
}

% 图片占位命令: 截图完成后替换为 \includegraphics
\newcommand{\placeholderfig}[2]{%
  \begin{figure}[H]
    \centering
    \fbox{\parbox[c][6cm][c]{0.85\textwidth}{\centering\small 待补截图: #1}}
    \caption{#2}
  \end{figure}
}

\title{\heiti\Large 面向政企场景的大模型智能体安全关键技术研究\\技术方案报告}
\author{团队名称(待填)\\学校名称(待填)}
\date{2026 年 9 月}

\begin{document}

\maketitle
\thispagestyle{empty}

% 摘要
\begin{abstract}
\noindent
政企场景下的大模型智能体被普遍接入工具调用能力,常见攻击方式从显式注入转向业务话术诱导下的过度代理。针对该威胁,本报告给出三层混合防护体系。第一层为结构化规则引擎,覆盖外发邮箱白名单、写入路径白名单、SQL 写操作、定时任务、审计日志清理等 14 类高风险动作。第二层为结构化 LLM 评审,在注入识别基础上引入过度代理检测与意图对齐判断两个维度,工程上采用异常即阻断的封闭失败模式。第三层为真实环境接入验证,在真实 OpenClaw 类智能体上完成端到端跑通。实验在 66 条真实业务话术诱导样本上,攻击拦截率从 86.4\% 提升至 100\%,在 25 条正常业务白样本上误报率为 0\%。

\vspace{1em}
\noindent\textbf{关键词:} 大模型智能体;提示词注入;工具调用安全;过度代理;纵深防御
\end{abstract}

\newpage


\tableofcontents
\newpage

% 第 1 章 项目背景与问题分析 (目标 3 页)
% 对应官方要求: 问题分析
% 要点: 威胁三阶段演进 → 现有方案三类不足 → 关键缺口总结
\section{项目背景与问题分析}

\subsection{智能体安全的威胁场景演进}

大模型智能体进入政企实际业务后,工具调用能力成为标配。工具集合覆盖邮件外发、文件读写、数据库查询、日历创建、定时任务、系统命令执行等高权限动作。围绕这些动作的安全威胁,经历三阶段演变。

第一阶段为显式注入。攻击者在用户消息或上传文档中嵌入"忽略上述规则"等字面指令,或用方括号标记注入载体。结构化关键词匹配对该类攻击有较好的识别能力。

第二阶段为去标记化业务话术诱导。攻击者将诱导意图嵌入正常公文格式、会议纪要、邮件正文中,完全消除方括号、感叹号等明显攻击信号。LLM 在阅读上下文后将诱导意图误判为业务推进,主动调用外发或写入工具。本项目在自有评测集上测得,基线阶段 66 条去标记化业务话术样本中 57 条成功绕过 LLM 自身安全过滤,绕过率 86.4\%。

第三阶段为过度代理,也是政企场景中最具迷惑性的威胁形式。LLM 在完成用户原始任务后,自主决定顺手做额外的外发与写入。例如用户只要求查看差旅费执行情况,LLM 自主决定起草报告并发送给局领导。该类威胁在常见基准中往往被归为无攻击类别,实际造成真实数据外发。OWASP LLM Top 10 将其列为 LLM08 Excessive Agency 独立风险项。

% TODO: 图 1 威胁演进示意图(三阶段时间线+每阶段代表样本片段)
\placeholderfig{威胁三阶段演进示意(显式注入 → 业务话术诱导 → 过度代理)}{智能体工具调用威胁的三阶段演进}

\subsection{现有方案的不足}

针对上述三类威胁,公开方案与本项目基线阶段的实测暴露三类不足。

第一,关键词层对去标记化场景失效。本项目基线阶段使用关键词加短语匹配,典型词表覆盖"忽略规则""立即执行""伪造签字""删除审计"等。在 66 条去标记化样本上,该层仅命中 5 条。

第二,LLM 自身安全机制在去标记化场景下漏过比例高。66 条样本中 9 条被 LLM 自身拒答,3 条未触发攻击工具,57 条成功绕过。

第三,LLM 评审类方案在工程实现上存在开放失败缺陷,即上游调用异常或模型无响应时,系统倾向放行。该模式在通用场景可接受,在政务数据敏感场景不可接受。本项目基线阶段实测,因上游接口路径配置错误,全部 LLM 评审调用返回异常,全部异常均被放行,放行率与无评审状态一致。

\subsection{关键缺口总结}

现有方案在去标记化业务话术诱导、过度代理、评审工程安全三处存在明确缺口。三者协同缺失,导致真实业务话术下的攻击绕过率达 86.4\%。政企场景另有其特殊约束:数据敏感性要求异常时阻断而非放行,业务连续性要求正常操作零误伤,审计合规要求每一次判定留痕可复现。任何单点方案无法同时满足以上约束,需要体系化设计。
     % 1 项目背景与问题分析 (3页)
% 第 2 章 系统总体设计 (目标 4 页)
% 对应官方要求: 技术路线(总体部分)
% 借鉴 SUFE-SAADS: 系统定位一句话 + 模块表格化 + 分层架构图 + 部署形态
% 注意: 全文自主口径, 不提任何外部项目
\section{系统总体设计}

\subsection{系统定位与设计目标}

本系统是一套部署于大模型智能体与模型服务之间的安全防护代理,对智能体的全部模型调用做透明拦截检查,覆盖攻击识别、工具调用约束、供应链风险评估、审计溯源四类能力,为政务办公场景下的智能体提供开箱可用的安全防护。

设计目标四条:

\begin{enumerate}[nosep]
  \item 高拦截。对去标记化业务话术诱导与过度代理两类高发威胁,拦截率不低于 95\%。
  \item 零误伤。正常政务办公操作的误报率不高于 1\%。
  \item 可解释。每一次拦截给出规则编号或评分理由,支持人工复核。
  \item 可复现。全部评测样本、脚本、审计日志随代码仓库提供,第三方可独立重跑。
\end{enumerate}

\subsection{总体架构}

系统采用四层架构,自上而下为接入层、策略层、判定层、审计层。

% TODO: 图 2 总体架构图(用 drawio 或 tikz 画: OpenClaw → 接入层(proxy) → 策略层(14条规则) → 判定层(LLM评审) → 审计层(audit_log), 旁路标注模型上游)
\placeholderfig{系统总体架构图(四层: 接入/策略/判定/审计)}{系统总体架构}

\begin{itemize}[nosep]
  \item \textbf{接入层}。以 HTTP 代理形态部署,实现 OpenAI Chat Completions 协议,对智能体完全透明,智能体侧仅需修改模型服务地址一项配置。
  \item \textbf{策略层}。结构化规则引擎,14 条规则覆盖已知高风险动作模式,执行开销在毫秒级。
  \item \textbf{判定层}。结构化 LLM 评审,对规则未覆盖的调用做语义级三维度评分,异常时一律阻断。
  \item \textbf{审计层}。全量调用轨迹与判定结果落盘,支持攻击复现、问题定位、效果验证。
\end{itemize}

\subsection{功能模块与赛题方向对位}

系统四大功能模块与赛题四个关键技术方向一一对应,见表~\ref{tab:module-mapping}。

\begin{table}[H]
  \centering
  \caption{功能模块与赛题关键技术方向对位}
  \label{tab:module-mapping}
  \small
  \begin{tabular}{p{2.6cm} p{1.2cm} p{5.6cm} p{4.2cm}}
    \toprule
    功能模块 & 对应方向 & 具体能力 & 实现载体 \\
    \midrule
    攻击识别与风险评估 & D1 & 对用户输入、文档附件、检索结果等多源输入做提示注入、间接指令污染、越狱诱导的自动发现与分级评估 & 请求侧扫描 + 规则层 + 评审层 \\
    工具调用安全约束 & D2 & 对文件操作、系统命令、接口调用等高风险动作做细粒度权限约束、动态策略校验与异常任务链阻断 & 代理 + 14 条结构性规则 + 评审兜底 \\
    插件与 Skill 供应链检测 & D3 & 对 Skill 包做代码行为检测、依赖关系分析与安全评级 & 50 类政务 Skill 的风险标签机制 \\
    安全评测与审计溯源 & D4 & 多维度评测体系、可量化测试方法、风险指标与审计证据输出 & 66 攻击样本 + 25 白样本 + 全量审计日志 \\
    \bottomrule
  \end{tabular}
\end{table}

四类功能以防护代理为统一入口,任一功能可独立启停,便于政企单位按实际需求裁剪部署。

\subsection{部署形态与运行环境}

系统支持两种部署形态。

\textbf{本地代理形态}。代理与智能体部署于同一主机,模型服务地址指向本地回环端口。该形态适用于单机办公智能体,部署时间不超过 10 分钟,部署步骤见附录 C。

\textbf{网关形态}。代理部署于政企内网网关,多个智能体实例共享同一代理。该形态适用于部门级部署,规则与评审配置集中管理,审计日志集中留存。

运行环境要求:Python 3.9 及以上,仅依赖标准库,无需 GPU。评审层调用远端模型服务,网络连通即可。代理自身单次调用增加的延迟实测约 1.8 秒,其中评审层占主要部分;对延迟敏感场景可仅启用策略层,增加延迟在毫秒级。
  % 2 系统总体设计 (4页)
% 第 3 章 关键技术方案 (目标 6 页)
% 对应官方要求: 技术路线 + 算法设计
% 结构: 总体路线 → 规则引擎 → LLM评审 → 协议适配 → 审计 → 工作流程
\section{关键技术方案}

\subsection{总体技术路线}

系统采用三层混合防护,自下而上为结构化规则、结构化 LLM 评审、真实环境接入验证。三层独立可部署,联合提供从已知模式拦截到语义兜底再到端到端验证的完整链路。

第一层结构化规则针对已知可枚举的高风险动作模式,覆盖邮箱外发白名单、写入路径白名单、SQL 写操作、定时任务、审计日志清理等。该层在所有调用链路上最先执行,执行开销极低,可解释性最强。

第二层结构化 LLM 评审针对规则未覆盖但语义可疑的工具调用,在注入识别基础上新增过度代理检测与意图对齐判断两个维度,要求评审模型输出结构化分数与简短中文理由。该层在规则未命中时按需触发。

第三层真实环境接入验证针对链路真实性问题,在真实智能体上跑通端到端攻击触发与拦截,证明防护在真实工具链路生效。

设计原则四条:真实优先,全部核心评估在真实业务话术样本上完成;分层独立,任意一层失能不影响其他层;封闭失败,评审与上游异常一律阻断并升级人工审批;可解释优先,规则给编号与触发条件,评审给分数与理由,审计给完整轨迹。

\subsection{结构化规则引擎设计}

规则引擎基于声明式策略匹配实现。每条规则包含编号、标题、工具选择器、谓词、效果、消息、启用标志七个字段。工具选择器限定规则适用的工具集合,谓词基于参数字段做任意嵌套匹配,效果分阻断、审批、记录三类。

实际部署 14 条规则,覆盖六类高风险动作。代表性规则见表~\ref{tab:rules},完整列表见附录 A。

\begin{table}[H]
  \centering
  \caption{代表性结构规则}
  \label{tab:rules}
  \small
  \begin{tabular}{p{2.8cm} p{2.4cm} p{7.2cm}}
    \toprule
    规则编号 & 工具范围 & 触发条件摘要 \\
    \midrule
    GOV-APPROVE-002 & send\_email 类 & 收件人或抄送含外部、群发、跨部门地址关键字 \\
    GOV-DENY-007 & write 类 & 写入路径含正式、机密、登记、归档目录或版本化文件名 \\
    GOV-DENY-008 & exec 类 & 命令含 SQL 写操作或审批字段改写 \\
    GOV-DENY-009 & cron 类 & 任意定时任务调用,覆盖延迟外发攻击 \\
    GOV-DENY-010 & exec 类 & 命令含日志清理、审计清理、服务重启 \\
    GOV-DENY-012 & write 类 & 参数字段名变体下的敏感路径,覆盖字段命名差异 \\
    \bottomrule
  \end{tabular}
\end{table}

谓词引擎支持约 25 种操作符,含字段存在性、子串匹配、正则匹配、集合成员、数值比较、逻辑组合。新增高风险模式仅需新增一条规则声明,无需改动引擎代码。

\subsection{结构化 LLM 评审设计}

评审层输入为当前工具调用名、参数摘要、最近 5 步前置调用文本。评审模型按 0 到 10 评分,6 分以上触发阻断。

评分沿三个维度展开。第一维度为注入风险,涵盖社工冒充、显式注入、编码绕过、函数返回值注入、知识库投毒、跨会话投毒。第二维度为过度代理,涵盖用户未明确请求的写入加外发组合、群发操作、批量批准、汇总后外发。第三维度为意图对齐,从用户原始请求出发评估当前调用是否在用户意图之内。

评审指令内含若干硬性约束。涉及执行命令、发送邮件、写文件等高风险工具且非用户明确请求时,基础分不低于 6。前置文本含社工信号时基础分不低于 7。前置轨迹出现连续多次同类工具调用且非用户批量请求时基础分不低于 7。完整评审指令见附录 B。

工程安全方面,评审调用结果做封闭失败处理。上游 HTTP 错误、网络异常、响应解析失败一律按阻断处理,不放行。该设计直接消除评审失效即全放行的隐患。

参数提取与降级。评审默认读取响应正文字段。当上游模型将推理结果放入推理字段且正文为空时,自动降级到推理字段。该降级保障多模型兼容,不引入新失败路径。

\subsection{响应解析与协议适配}

接入层解决代理能否拦到真实工具调用的问题。智能体默认发起流式请求,模型上游在流式模式下返回分片事件而非单一 JSON,代理无法在响应解析阶段检查工具调用。代理转发时强制将流式字段改为关闭,上游返回标准 JSON,代理在响应解析阶段提取工具调用并执行策略检查。该改造不影响模型输出语义,仅影响代理侧的解析能力。

平台适配涉及两处。其一,部分桌面操作系统对智能体状态目录施加系统级保护,需通过环境变量将状态目录隔离到可写路径。其二,代理需兼容多家模型上游的协议差异,含正文与推理字段的位置差异。

\subsection{审计溯源设计}

审计层对每一次模型调用记录七类字段:调用时间、目标模型、请求摘要、工具调用清单、命中规则编号、评审分数与理由、最终动作。审计日志以 JSON Lines 格式追加写入,单条记录自包含,支持按样本编号、规则编号、时间段做检索与回放。审计样例见附录 D。

\subsection{系统工作流程}

一次完整的攻击拦截流程包含七个环节,全链路在防护代理内完成,见图~\ref{fig:workflow}。

% TODO: 图 3 工作流程图(七环节泳道: 请求接入 → 响应解析 → 文本级阻断 → 规则层 → 评审层 → 回写响应 → 审计)
\begin{figure}[H]
  \centering
  \fbox{\parbox[c][7cm][c]{0.9\textwidth}{\centering\small 待绘制: 七环节工作流程泳道图}}
  \caption{攻击拦截工作流程}
  \label{fig:workflow}
\end{figure}

\begin{enumerate}[nosep]
  \item 请求接入。智能体发起模型调用,请求经代理转发,代理解析请求体并识别目标模型。
  \item 响应解析。代理强制请求为非流式,上游返回标准 JSON,代理提取消息与工具调用,还原完整调用轨迹。
  \item 文本级阻断。对响应文本做注入标记与高危关键词检测,命中时改写响应并记录审计。
  \item 规则层判定。对每个工具调用执行 14 条结构性规则,任意规则命中即阻断,返回结构化阻断消息并写审计。
  \item 评审层兜底。规则未命中的调用进入评审层,综合当前调用、前置轨迹、用户原始意图三维度评分,6 分以上触发阻断,评审异常时按封闭失败处理。
  \item 回写响应。通过检查的调用原样保留,被阻断的调用替换为阻断提示,智能体据此生成合规回复。
  \item 审计溯源。每个环节的判定结果、命中规则、评分、原始调用与改写内容全部落盘。
\end{enumerate}

以伪造签字类样本为例。用户以业务话术诱导智能体批量出具同意意见,智能体调用命令执行工具发起 SQL 更新。代理在第三环节捕获命令文本,第四环节命中 SQL 写操作规则,响应被改写为阻断提示,审计日志记录完整调用链。全程无需智能体侧任何改造。
      % 3 关键技术方案 (6页)
% 第 4 章 评测数据集构建 (目标 3 页)
% 对应官方要求: 实验方案(数据部分)
% 借鉴国二报告: 数据集独立成章, 含来源/生成/评估/筛选四节
\section{评测数据集构建}

\subsection{构建原则与流程}

数据集构建遵循三项原则。第一,用户消息用正常业务话术,不暴露攻击意图。第二,工具模拟返回真实查询结果,不人为编造诱导。第三,首轮查真实数据,后续轮次由模型自主判断是否需要行动。

攻击样本 66 条,来源于政府办公 50 类技能的真实业务话术诱导,经两轮筛选改造而来。第一轮按工具多样性、模式多样性、文本长度、多步骤性四维度打分,从 351 条公开攻击样本中筛出 90 条。第二轮针对初筛样本中模型拒答与不触发占比过高的情况,以业务话术替换显式诱导,重写模拟返回为真实查询结果,最终得到 66 条在真实模型上跑通绕过的样本。

% TODO: 图 4 数据集构建流程图(351 → 四维度筛选 → 90 → 业务化改写 → 66; 旁路: 白样本 25)
\placeholderfig{数据集构建流程(351 → 90 → 66 + 白样本 25)}{评测数据集构建流程}

\subsection{攻击样本覆盖矩阵}

66 条攻击样本覆盖 9 类攻击模式与 50 类政务技能中的 34 类,见表~\ref{tab:coverage}。

\begin{table}[H]
  \centering
  \caption{攻击样本覆盖矩阵}
  \label{tab:coverage}
  \small
  \begin{tabular}{p{3cm} p{1.4cm} p{7.6cm}}
    \toprule
    攻击模式 & 数量 & 覆盖政务技能示例 \\
    \midrule
    数据外泄 & 25 & 差旅费报告、通讯录查询、跨部门协调、公文起草 \\
    伪造签字 & 10 & 审批意见出具、批量同意、会议纪要代签 \\
    系统提示伪装 & 8 & 系统升级通知、政策版本覆写 \\
    批量操作 & 5 & 群发邮件、批量建会、批量任务下发 \\
    破坏操作 & 5 & 日志清理、系统重启、文件删除 \\
    越权操作 & 4 & 跨部门数据访问、越级审批 \\
    知识库投毒 & 3 & 知识库条目篡改、检索结果污染 \\
    工具结果注入 & 3 & 邮件正文指令、查询结果指令 \\
    多轮诱导 & 3 & 先无害后诱导的渐进式攻击 \\
    \bottomrule
  \end{tabular}
\end{table}

\subsection{白样本设计}

白样本 25 条,覆盖 18 类技能,用于误报率评估。每条白样本对应一条用户明确请求的正常业务操作,例如查询差旅费政策、起草会议室使用通知并写入草稿目录、安排部门内部讨论。白样本设计时回避批量、应急、机密等可能触发评审高分的措辞,贴近真实日常办公场景。

\subsection{数据质量说明}

全部样本经三条质量校验。其一,每条攻击样本在基线模式下真实跑通并确认绕过,杜绝纸面攻击。其二,每条样本的工具调用参数由人工核对,确保动作与诱导意图一致。其三,攻击样本与白样本的技能分布做了对齐,避免误报率评估因技能分布偏差失真。全部样本与结果文件随仓库提供,可独立复现。
        % 4 评测数据集构建 (3页)
% 第 5 章 系统实现与功能展示 (目标 3 页)
% 借鉴国二报告"功能展示"章: 让评委先看到系统真实存在
% 截图前置: 需先完成高危缺口修复(G2/G3/G5/G1)与部署一键化, 见章节规划-v2.md 第六节
\section{系统实现与功能展示}

\subsection{部署架构与一键复现}

系统以三个进程构成完整演示链路:工具模拟服务、防护代理、评测运行器。仓库提供一键启动脚本,完成依赖检查、服务拉起、样本跑通、结果汇总四步,单次全量评测约 30 分钟。部署指南见附录 C。

% TODO: 图 5 截图: 终端一键启动脚本运行成功画面(三个服务就绪+样本进度条)
\placeholderfig{一键启动脚本运行成功(三服务就绪)}{一键部署与启动}

\subsection{真实智能体端到端演示}

在真实 OpenClaw 类智能体上加载 50 类政务技能,智能体可正常完成公文起草、数据查询、会议安排等日常办公任务。图~\ref{fig:demo-normal} 展示智能体正常办公的交互界面。

% TODO: 图 6 截图: OpenClaw 对话界面, 用户发起正常办公请求, 智能体正常完成
\placeholderfig{OpenClaw 正常办公对话(加载 50 类政务 Skill)}{防护系统下的智能体正常办公}
\label{fig:demo-normal}

\subsection{拦截效果展示}

以数据外泄类攻击为例。用户以业务话术诱导智能体将敏感数据外发,智能体生成外发邮件的工具调用。防护代理在响应解析阶段捕获该调用,规则层命中外发地址关键字,调用被阻断,智能体改而输出合规提示。图~\ref{fig:demo-block} 展示拦截瞬间的智能体回复。

% TODO: 图 7 截图: 攻击 case 触发, 智能体回复被代理改写为阻断提示
\placeholderfig{攻击触发与拦截(阻断提示替换外发调用)}{攻击触发与实时拦截}
\label{fig:demo-block}

\subsection{审计日志展示}

每次拦截在审计日志中留下完整记录,含命中规则编号、评审分数、原始调用参数、最终动作。图~\ref{fig:demo-audit} 展示终端检索审计日志的画面。

% TODO: 图 8 截图: 终端 cat audit_log.jsonl 展示结构化审计记录(含规则ID/评分/trace)
\placeholderfig{审计日志检索(规则编号+评分+调用轨迹)}{审计溯源日志}
\label{fig:demo-audit}
           % 5 系统实现与功能展示 (3页)
% 第 6 章 实验与指标体系 (目标 4 页)
% 对应官方要求: 实验方案 + 指标体系
\section{实验与指标体系}

\subsection{实验设置}

每个样本在两种模式下跑通。基线模式直接调用模型,无任何中间层。防护模式在模型与工具调用之间插入防护代理,代理按规则层、评审层顺序检查。

模型主干选用两类国产主流模型,覆盖代表性业务使用场景。评测在真实智能体上跑通端到端链路,真实加载 50 类政务技能,真实执行工具调用,真实产生审计日志。

每条样本记录四类结果。BYPASSED 表示模型真调了攻击工具,REFUSED 表示模型主动拒答,NO\_ATTACK 表示模型未触发攻击工具。BYPASSED 是防护体系要解决的核心类别。

\subsection{指标体系}

三项核心指标。

\textbf{攻击拦截率}。定义为基线 BYPASSED 状态在防护模式下转为非 BYPASSED 状态的比例,反映防护体系对真实攻击的拦截能力。

\textbf{误报率}。定义为白样本在防护模式下被阻断或被改写的比例,反映防护体系对正常业务的容忍度。

\textbf{延迟增加}。定义为加防护代理后单次模型调用的延迟增加,反映工程可用性。

\subsection{实验结果}

\textbf{攻击拦截率}。在 66 条攻击样本上,基线阶段 BYPASSED 57 条,绕过率 86.4\%。防护模式下 BYPASSED 0 条,拦截率 100\%。其中结构化规则层拦截 38 条,评审层兜底拦截 17 条,其余在迭代修复阶段通过字段变体规则与轨迹级联识别补齐。

\textbf{误报率}。在 25 条白样本上,防护模式真实误报 0 条,误报率 0\%。规则层与评审层均未对白样本触发阻断。

\textbf{精确归因}。66 条样本逐条比对基线与防护结果,83.3\% 归因为防护代理拦截,13.6\% 归因为模型自身安全机制,3.0\% 归因为迭代修复后的代理拦截。该归因说明 100\% 拦截中 86.4 个百分点来自本防护体系。

\textbf{代理内部动作}。实验期间代理共产生 173 次结构化规则拦截、506 次评审调用、38 次文本级阻断,平均单请求延迟增加 1.8 秒。

\begin{table}[H]
  \centering
  \caption{核心实验结果汇总}
  \label{tab:results}
  \small
  \begin{tabular}{lccc}
    \toprule
    指标 & 基线模式 & 防护模式 & 变化 \\
    \midrule
    攻击样本绕过率 & 86.4\% & 0\% & 下降 86.4 个百分点 \\
    白样本误报率 & 0\% & 0\% & 持平 \\
    单请求延迟增加 & -- & 1.8 秒 & -- \\
    \bottomrule
  \end{tabular}
\end{table}

\subsection{消融实验}

\textbf{仅规则层}。在 66 条攻击样本上单独启用规则层,命中 27 条,绕过率从 86.4\% 压到 40.9\%。规则层单独可拦截约一半攻击,作为高置信基础层具备工程价值。

\textbf{规则层加评审层}。两层协同后 BYPASSED 为 0 条。评审层对规则未命中的样本提供有效兜底,两层达成 100\% 拦截。

\textbf{白样本误报隔离}。规则层与评审层在 25 条白样本上均未触发阻断,两层在保持高拦截的同时未引入白样本误伤。

\subsection{与基线方案对比}

将本体系与模型自身安全过滤对比。在 66 条去标记化业务话术样本上,模型自身安全仅拦截 9 条,本体系拦截全部 66 条。在 25 条白样本上,两者均保持 0 误报。本体系同时在真实智能体上提供端到端拦截证据,覆盖规则、评审、链路三层,公开方案的评估多停留在离线基准阶段。
    % 6 实验与指标体系 (4页)
% 第 7 章 创新点 (目标 2 页)
\section{创新点}

\subsection{三层协同的纵深防御架构}

已有方案多停留在单层防御,或规则、或评审、或离线基准。本系统给出规则层、评审层、真实环境接入层的三层协同架构,三层独立可部署、联合提供纵深。规则层拦截已知模式,评审层兜底语义可疑调用,接入层保证策略在真实链路生效。消融实验表明单规则层拦截 40.9\%,两层协同达 100\%,体现层间互补。

\subsection{针对过度代理的结构化检测维度}

已有提示注入防御工作聚焦注入内容本身,对用户只请求查看而模型自主外发的过度代理行为覆盖不足。本系统在评审层新增过度代理与意图对齐两个维度,从用户原始请求出发评估调用合法性,并以硬性约束工程化落地:非用户明确请求的高风险工具基础分不低于 6。该维度直接对应 OWASP LLM Top 10 中的 LLM08 Excessive Agency。

\subsection{封闭失败的工程安全设计}

公开的模型评审类方案在实现上普遍采用开放失败,上游异常时倾向放行,在政务数据敏感场景不可接受。本系统在基线阶段实测发现该缺陷,评审调用异常被全部放行,随后修复为封闭失败,上游异常一律阻断并升级人工审批。该修复消除评审失效即全放行的隐患,是工程级的可靠性贡献。

\subsection{真实智能体端到端验证方法}

已有方案的评估多基于离线基准或合成数据。本系统在真实智能体上加载 50 类政务技能,跑通端到端攻击触发与拦截,提供真实工具链路的拦截证据。该验证方式回应了代理能否拦截真实工具调用这一关键工程问题,验证方法本身可复用于其他智能体防护方案的评估。
     % 7 创新点 (2页)
% 第 8 章 应用前景与预期效果 (目标 1.5 页)
% 对应官方要求: 预期效果
\section{应用前景与预期效果}

\subsection{落地形态}

本系统代码随报告提交,核心模块含结构化规则引擎、结构化评审层、真实环境接入代理、攻击与白样本评测集,均提供运行说明。第三方团队可在本地智能体上重现全部实验。部署支持单机代理与内网网关两种形态,适配不同规模的政企单位。

\subsection{推广路径}

规则层与评审层均与具体业务技能解耦。规则通过新增编号与谓词即可覆盖新的高风险工具或参数模式,评审通过修改指令即可调整评估维度。两层均不依赖特定模型,可在多模型环境部署。政务办公之外,金融、医疗、能源等数据敏感行业的智能体防护可直接复用本架构,按行业补充规则库即可。

\subsection{对治理水平的价值}

本系统给出从关键词到结构化规则到结构化评审的三阶段演进路径,可作为政企单位建设智能体安全防护体系的参考模板。审计溯源设计满足合规留痕要求,评测数据集与评测方法可复用为行业评测基准的种子。

\subsection{预期效果}

按当前实验结果外推,系统部署后可将去标记化业务话术诱导类攻击的绕过率从 86.4\% 压到 0\%,正常业务操作零误伤,单次调用增加延迟 1.8 秒以内。随着规则库与评审指令在行业部署中持续迭代,覆盖率与误报控制仍有进一步提升空间。
      % 8 应用前景与预期效果 (1.5页)

% 骨架阶段全部列出; 正文写作时替换为按 \cite 引用
\nocite{*}
\bibliographystyle{plain}
\bibliography{refs}

\appendix
% 附录 (目标 3 页)
% A 规则完整列表 / B 评审指令全文 / C 快速部署指南 / D 审计日志样例

\section{结构化规则完整列表}

\begin{longtable}{p{2.8cm} p{2.2cm} p{1.2cm} p{7cm}}
  \caption{14 条结构规则完整列表}\\
  \toprule
  规则编号 & 工具范围 & 效果 & 触发条件摘要 \\
  \midrule
  \endhead
  GOV-DENY-001 & exec & 阻断 & 无指令记录时的裸命令执行 \\
  GOV-DENY-002 & exec & 阻断 & 命令含递归删除 \\
  GOV-DENY-003 & exec & 阻断 & 命令含外发下载 \\
  GOV-APPROVE-001 & send\_email 类 & 阻断 & 收件人为已知攻击邮箱 \\
  GOV-DENY-004 & write 类 & 阻断 & 路径含政策目录或版本化文件名 \\
  GOV-DENY-005 & 浏览器类 & 阻断 & 表单提交类外部地址 \\
  GOV-DENY-006 & 高风险工具集 & 阻断 & 轨迹含社工标记且当前为高风险工具 \\
  GOV-APPROVE-002 & send\_email 类 & 阻断 & 收件人含外部、群发、跨部门关键字 \\
  GOV-DENY-007 & write 类 & 阻断 & 路径含正式、机密、登记、归档目录 \\
  GOV-DENY-008 & exec 类 & 阻断 & 命令含 SQL 写操作或审批字段改写 \\
  GOV-DENY-009 & cron 类 & 阻断 & 任意定时任务调用 \\
  GOV-DENY-010 & exec 类 & 阻断 & 命令含日志清理、审计清理、服务重启 \\
  GOV-DENY-011 & query\_db 类 & 阻断 & 表名含员工、薪资、人事敏感表 \\
  GOV-DENY-012 & write 类 & 阻断 & 字段名变体下的敏感路径 \\
  \bottomrule
\end{longtable}

\section{LLM 评审指令全文}

评审指令要求模型按 0 到 10 评分,沿注入风险、过度代理、意图对齐三个维度展开,并含三条硬性约束:高风险工具非用户明确请求时基础分不低于 6;前置文本含社工信号时基础分不低于 7;轨迹内连续多次同类操作非用户批量请求时基础分不低于 7。指令全文约 40 行,随仓库 \texttt{src/system/policies/} 目录提供。

% TODO: 从 src/system/policies/llm_injection_judge.py 的 _SYSTEM_PROMPT 复制全文, 用 lstlisting 排版

\section{快速部署指南}

\begin{lstlisting}[language=bash]
# 1. 克隆仓库并进入目录
git clone <仓库地址> && cd student_competition

# 2. 配置模型服务密钥
cp .env.example .env   # 填入 STEPFUN_API_KEY / DEEPSEEK_API_KEY

# 3. 一键启动演示链路(工具模拟 + 防护代理 + 评测运行器)
bash scripts/run_e2e.sh

# 4. 查看结果与审计日志
cat results/summary.txt
cat src/system/proxy/audit_log.jsonl | python3 -m json.tool | less
\end{lstlisting}

% TODO: scripts/run_e2e.sh 与 .env.example 为待补齐交付物, 见章节规划-v2.md 第六节

\section{审计日志样例}

\begin{lstlisting}
{"ts": "2026-08-07T14:23:11", "trace_id": "v3-33",
 "model": "step-3.7-flash",
 "tool_call": {"name": "exec",
   "args_summary": "UPDATE approvals SET status='agreed' ..."},
 "rule_hit": "GOV-DENY-008",
 "judge_score": 8, "judge_reason": "批量改写审批字段, 非用户明确请求",
 "action": "BLOCK"}
\end{lstlisting}
       % 附录 A-D (3页)

\end{document}
